% Online Resource 2 for:
%   E. H. Brookes, "Exactness of an iterated divide-and-conquer method
%   for large-scale nonnegative least squares", Optimization Letters.
% Springer publishes supplementary PDFs as supplied; typeset to match
% the article (sn-jnl class, sn-mathphys-num style).
\documentclass[pdflatex,sn-mathphys-num]{sn-jnl}

\usepackage{amsmath,amssymb,amsfonts}

\raggedbottom

\hypersetup{
  pdftitle={Online Resource 2: Notes on the UltraScan implementation},
  pdfauthor={Emre H. Brookes},
}

\begin{document}

\title[Online Resource 2]{Online Resource 2:\\ Notes on the UltraScan
implementation}

\author*[1]{\fnm{Emre H.}
\sur{Brookes}\,\orcid{https://orcid.org/0000-0001-7348-6476}}%
\email{emre.brookes@umontana.edu}

\affil[1]{\orgdiv{Department of Chemistry \& Biochemistry},
\orgname{University of Montana}, \orgaddress{\city{Missoula},
\state{Montana}, \country{USA}}}

\maketitle

\noindent
Supplementary material for \emph{Exactness of an iterated
divide-and-conquer method for large-scale nonnegative least squares}
(Optimization Letters). This note records how the divide-and-conquer
NNLS implementation shipped in the open-source UltraScan package
relates to the article's conditions (H1)--(H3) and to its stopping
certificate.
Source locations refer to the repository
\texttt{https://github.com/ehb54/ultrascan3} at tag v4.1.0 (commit
\texttt{c31386b4a}); numbered results are the article's. Sections
2--3 describe the noise-free configuration
(\texttt{noisflag}~$=0$): when time- or radially-invariant noise
fitting is enabled, the block solver takes a different code path
whose relationship to (H1) is under separate investigation, and the
transfer statements below are not asserted for that configuration.

\section*{1\quad Merge and feedback structure: (H2) and (H3)}

The MPI 2DSA master
(\texttt{programs/us\_mpi\_analysis/2dsa\_master.cpp}, functions
\texttt{process\_solutes} and \texttt{iterate}) accumulates each
worker's complete solute vector into a depth-indexed bin, ships a bin
whole to the next depth when adding the following result would exceed
the size cap, and, at each refinement iteration, unions the previous
iteration's nonzero solutes into every subgrid job with duplicates
removed. Returned solute sets are therefore never split across merge
jobs, and every leaf carries the incumbent --- the canonical
instantiation of conditions (H2) and (H3), in the atomic
whole-set-packing form the article's Section~8 recommends, with bins
closed by an adaptive size cap rather than a fixed fan-in. The
interactive 2DSA program (\texttt{programs/us\_2dsa}), a separate
threaded implementation, does the same.

\section*{2\quad Grid screening: the scope of (H1)}

The block solver screens the grid before solving: every column whose
Euclidean norm falls below a cutoff $c$ is dropped from the
subproblem (\texttt{utils/us\_solve\_sim.cpp};
$c=\texttt{\_NORM\_CUTOFF\_}=1.00$ unless overridden for the run
through debug text). The screen depends only on the
column, not on the block or the round, so the deployed method is
Algorithms~1--2 of the article applied to the screened set
$S'=\{j:\lVert a_j\rVert_2\ge c\}$; every result, including
Theorem~1's certificate, transfers verbatim with $S'$ in place of
$S$, and $R\subseteq S'$ holds automatically since a screened column
never carries a positive amplitude. The scope consequence is that the
certificate then asserts global optimality over the \emph{screened}
grid rather than the nominal one; the screen is a heuristic discard of
the kind the article's Section~1 contrasts with safe screening --- a
column with $\lVert a_j\rVert_2$ just below the cutoff can still price
positively against the merged residual.

\section*{3\quad Stopping logic: which exits are certified}

The refinement loop of the MPI master has four exits, checked after
each pass in this order. \emph{Cap:} the pass count is compared with
the iteration cap before anything else, so a stopping test at round
$k$ can fire only if the cap exceeds $k$; the cap comes from the job
request (\texttt{parameters["max\_iterations"]}), is raised from 1 to 5 in
single-master Monte Carlo runs; parallel-masters Monte Carlo runs do
not refine at all.
\emph{Rollback:} past half the cap, a pass whose variance worsens ends
refinement. \emph{Value:} the variance change falls below $10^{-100}$
(\texttt{min\_variance\_improvement}). \emph{Solute set:} the solute
set repeats with all concentrations within 10\%, a test stricter than
support repetition. The solute-set exit is certified by
Theorem~1(b); the cap and the rollback never are. The value exit ---
exact value equality in effect --- is certified too, provided no data
point reaches the optical-density limit: such points are zeroed, with
their rows of the design matrix, in the problem solved but still
counted in the reported variance, which then no longer tracks the
objective. Under the same proviso the article's Lemma~6 rules the
rollback out in exact arithmetic. The certified tests fired at rounds
4--13 in the article's correlated-kernel suites, so a cap of 4 or less
would bind first in every such trial.

The interactive program stops when its solute set repeats --- the
article's support test, certified by Theorem~1(b) --- or when the
per-point variance changes by at most an absolute $10^{-12}$, a
relaxation that differs from the relative test of the article's
Section~7 and is likewise not covered by Theorem~1(b). Its iteration
count (\texttt{us\_analysis\_control\_2d.cpp}) defaults to 1, which
means refinement is off: the default interactive run is the
non-iterated scheme of the article's Proposition~1. Enabling
refinement sets the count to 3, permitting one comparison, of round 2
with round 1.

\smallskip
\noindent The practical rule these notes support: treat a run as
certified by the article's theory exactly when it ends on a certified
exit --- on the MPI master, the solute-set exit, or the value exit when
no data point reaches the optical-density limit (the job log names the
exit that fired); in the interactive program, the solute-set test,
which it does not report separately --- and read the
certificate as global optimality over the screened grid. None of the
above bears on the article's theorems; it locates the shipped code
relative to their hypotheses.

\end{document}
